\documentclass[book.tex]{subfiles}

\begin{document}
\chapter{Electric Current and the Magnetic Field}
\label{chap:electric_current}
\noindent\rule{11cm}{0.4pt} \\
\textbf{Prerequisites:} \autoref{chap:electrostatics},  \\
\textbf{Difficulty Level:} * \\
\noindent\rule{11cm}{0.4pt}
\vspace{5mm}


In \autoref{chap:electrostatics}, we learned about the behavior of systems of point charges. In the real world, it's more common to deal with flowing electric charges or electric currents. In this chapter we introduce the mathematical and computational models used to study electric current. We also introduce the notion of the magnetic field, a powerful abstraction useful for the study of complex electric systems. Magnetic fields are associated with electric currents much as electric fields are associated with static electrical charges.

\section{Electric Current}

Let's define electric current more formally. We say that the electric current is the rate of flow of electric charge past a point or a region. The mathematical symbols used to define electrical current are usually either $I$ or $J$. Suppose that charge $Q$ flows through a given region of space in time $t$. We can then define
\begin{align}
    I = \frac{Q}{t}
\end{align}

As the time $t$ goes to $0$, we can replace this definition with a suitable limit
\begin{align}
    I = \frac{dQ}{dt}
\end{align}
We can define types for these quantities
\begin{lstlisting}
Q : Ampere, t: sec, I: Ampere/sec
\end{lstlisting}
Recall that \lstinline{Ampere} is the same type as $\mathbb{R}$ and \lstinline{sec} is the same type as $\mathbb{R}^+$. Since $Q \in \mathbb{R}$ is a real valued quantity, and $t \in \mathbb{R}^+$ is a nonnegative quantity, we can say that $I \in \mathbb{R}$ is also a real valued quantity.

\begin{marginfigure}
% \includegraphics[width=\textwidth]{figures/Physics/electromagnetism/electric_current/electric_current_charges.png}
 \includegraphics[width=\textwidth]
{figures/Physics/electromagnetism/electric_current/Electric current.png}
\caption{The current is given by the rate flow of charge through a given region of space.}%\url{https://commons.wikimedia.org/wiki/File:Electric_charge_and_electric_current.svg}}
\end{marginfigure}

%\textcolor{red}{TODO: Add an example for the physical theory of an electric field created by current flowing through a wire. }

\section{Lorentz Force Law}

Empirically, Coulomb's law isn't a complete description of the physics moving electrical particles since it ignores magnetism. Consequently, we introduce a new law, the \textit{Lorentz force law}. Let $\vec{E}(t)$ be the electric field and let $\vec{B}(t)$ be the magnetic field. A particle with charge $q$ moving with velocity $\vec{v}$ experiences a force given by
\begin{align}
    \vec{F} = q(E + v \times B)
    \label{lorentz:force}
\end{align}

We've introduced a new concept in this definition, that of the magnetic field $\vec{B}$. What is the magnetic field? To give a more sophisticated derivation of the magnetic field, we will need to wait until the section on special relativity. For now, it's not wrong to view Equation \ref{lorentz:force} as \textit{defining} the magnetic field. That is, the magnetic field can be viewed as an empirical correction term that fixes Coulomb's law to work for electrodynamical systems. The magnetic field is a field much like the classical field and has a similar type
\begin{lstlisting}
B: $\mathbb{R}^3 \to T\mathbb{R}^3$
\end{lstlisting}


\subsection{A charged particle in a magnetic field} 

In this simple example, we consider a charged particle moving in a magnetic field as illustrated in Figure~\ref{fig:charged_particle_magnetic_field}.

\begin{figure}[ht]
   % \includegraphics[width=\linewidth]{figures/Physics/electromagnetism/electric_current/charged_particle_magnetic_field.png}
    \includegraphics[width=\linewidth]{figures/Physics/electromagnetism/electric_current/Charged Particle in mag field.png}
   \caption{A charged particle in a magnetic field. }
\label{fig:charged_particle_magnetic_field}
\end{figure}
% TODO: This figure is sourced from \url{http://web.mit.edu/sahughes/www/8.022/lec10.pdf}, section 10.2. Replace with our own illustration.
Let's introduce a coordinate basis of orthonormal vectors for this example. We are working in $\mathbb{R}^3$ and we have vectors $\hat{i}, \hat{j}, \hat{k}$. We choose this basis such that $\hat{i} \times \hat{j} = \hat{k}$. Let's say that the $\vec{v_0} = \hat{i}$ and $\vec{B} = \hat{j}$ everyone in $\mathbb{R}^3$. The states of the system are 
\begin{align}
\mathcal{S} = \mathbb{R}^3 \times \mathbb{R}^3 \times \mathbb{R}^3  
\end{align}
Each state consists of a tuple $(\vec{x}, \vec{v}, \vec{a})$ of position, velocity, and acceleration for the charged particle. We define force $F$ by  
\begin{align}
F(x) = q(\vec{v} \times \vec{B}) 
\end{align}

Using the known value of $\vec{v_0}$, we can compute
\begin{align}
\vec{F_0} &= q (\vec{v_0} \times \vec{B}) \\
&= q \|\vec{v_0}\| \|\vec{B}\|(\hat{i} \times \hat{j}) \\
&= q\|\vec{v_0}\| \|\vec{B}\|\hat{k}
\end{align}

For the domain of applicability, we limit the velocity to non-relativistic speeds $\|\vec{v}\| < C$.

The parameters of the system are 
\begin{align}
\mathcal{W} = \{q, \vec{B}, \vec{v_0}\}
\end{align}
where $q$ is the charge of the particle, $\vec{B}$ is the constant magnetic field, and $\vec{v_0}$ is the initial velocity. The hyperparameters are given by $\mathcal{H} = \{C\}$.


%\textcolor{red}{TODO: Add definition for centripetal force}

%\textcolor{red}{TODO: Add example for classical circular motion.}

Note that the initial force $\vec{F_0}$ we computed points perpendicular to the motion of the particle. This is exactly the force needed to induce circular motion. Note that this motion will ensure that the charged particle is always perpendicular to the magnetic field so this circle will close in on itself. 

We can solve for the radius of the circle using the definition of centripetal force.
\begin{align}
    q\|\vec{v_0}\|\|\vec{B}\| &= \frac{m\|v_0\|^2}{R} \\
    R &= \frac{m\|v_0\|}{q\|\vec{B}\|}
\end{align}

\section{Amp\`{e}re's Law}

Where does the magnetic field come from? The rough idea is that magnetic fields are induced by movement of electric charge. Let's introduce one more law to formalize this, \textit{Amp\`{e}re's Law}. This law states that the line integral of the magnetic field $\vec{B}$ around a close path is a multiple of the current enclosed by the path.
\begin{equation}
\oint_C \vec{B} \dot d\vec{s} = \mu_0 I_{\textrm{encl}}
\end{equation}

\begin{figure}[ht]
   % \includegraphics[width=0.7\linewidth]{figures/Physics/electromagnetism/electric_current/current_in_curve.png}
    \includegraphics[width=0.8\linewidth]{figures/Physics/electromagnetism/electric_current/Current in curve.png}
   \caption{A curve enclosing a wire with current. }
\label{fig:current:in:curve}
\end{figure}
%TODO: This figure is sourced from \url{http://web.mit.edu/sahughes/www/8.022/lec10.pdf}, section 10.4. Replace with our own illustration.


\end{document}
